% -*- coding: utf-8 -*-
% !TEX program = xelatex
\documentclass[plain,noanswer,medmath]{../../template/jnuexam}
\usepackage{../../template/kysx}

\begin{document}


\examtitle{2017}{3}


\loadproblems{1}{2017P1}%载入试卷一
\loadproblems{2}{2017P2}%载入试卷二


\makepart{选择题}{1～8小题，每小题4分，共32分}


\useproblem{1}{1}{1}%【同试卷一第一(1)题】


\begin{problem}
二元函数$z = xy(3 - x - y)$的极值点是\pickout{}
\begin{abcd}
\item $(0,0)$．
\item $(0,3)$．
\item $(3,0)$．
\item $(1,1)$．
\end{abcd}
\end{problem}


\useproblem{1}{1}{2}%【同试卷一第一(2)题】


\begin{problem}
设函数$\sum_{n = 2}^{\infty}\left[\sin\frac{1}{n} - k\ln\left(1 - \frac{1}{n}\right)\right]$收敛，
则$k =$\pickout{}
\begin{abcd}
\item $1$．
\item $2$．
\item $- 1$．
\item $- 2$．
\end{abcd}
\end{problem}


\useproblem{1}{1}{5}%【同试卷一第一(5)题】


\useproblem{1}{1}{6}%【同试卷一第一(6)题】


\begin{problem}
设$A,B,C$为三个随机事件，且$A$与$C$相互独立，$B$与$C$相互独立，
则$A \cup B$与$C$相互独立的充要条件是\pickout{}
\begin{abcd}
\item $A$与$B$相互独立．
\item $A$与$B$互不相容．
\item $AB$与$C$相互独立．
\item $AB$与$C$互不相容．
\end{abcd}
\end{problem}


\useproblem{1}{1}{8}%【同试卷一第一(8)题】


\makepart{填空题}{9～14小题，每小题4分，共24分}


\begin{problem}
$\int_{- \pi}^{\pi}(\sin^3 x + \sqrt{\pi^2 - x^2}) \dx =$ \fillin{}．
\end{problem}


\begin{problem}
差分方程$y_{t + 1} - 2y_t = 2^t$的通解为$y_t =$ \fillin{}．
\end{problem}


\begin{problem}
设生产某产品的平均成本$\overline{C} (Q) = 1 + \e^{- Q}$，其中$Q$为产量，则边际成本为 \fillin{}．
\end{problem}


\useproblem{2}{2}{12}%【同试卷二第二(12)题】


\useproblem{1}{2}{13}%【同试卷一第二(13)题】


\begin{problem}
设随机变量$X$的概率分布为$P\{X = - 2\} = \frac{1}{2}$，$P\{X = 1\} = a$，$P\{X = 3\} = b$，
若$EX = 0$，则$DX =$ \fillin{}．
\end{problem}


\makepart{解答题}{15～23小题，共94分}


\useproblem[points=10]{2}{3}{15}%【同试卷二第三(15)题】


\begin{problem}[points=10]
计算积分$\iint_D \frac{y^3}{(1 + x^2 + y^4)} \dx\dy$，
其中$D$是第一象限中以曲线$y = \sqrt{x}$与$x$轴为边界的无界区域．
\end{problem}


\useproblem[points=10]{1}{3}{16}%【同试卷一第三(16)题】


\begin{problem}[points=10]
已知方程$\frac{1}{\ln(1 + x)} - \frac{1}{x} = k$在区间$(0,1)$内有实根，求$k$的范围．
\end{problem}


\begin{problem}[points=10]
设$a_0 = 1$，$a_1 = 0$，$a_{n + 1} = \frac{1}{n + 1}(na_n + a_{n - 1})$（$n = 1,2, \cdots$），
$S(x)$为幂级数$\sum_{n = 0}^{\infty}a_n x^n$的和函数．
\begin{enumerate}
\item 证明幂级数$\sum_{n = 0}^{\infty}a_n x^n$的收敛半径不小于$1$；
\item 证明$(1 - x)S'(x) - xS(x) = 0$（$x \in(- 1,1)$），并求$S(x)$的表达式．
\end{enumerate}
\end{problem}


\useproblem[points=11]{1}{3}{20}%【同试卷一第三(20)题】


\useproblem[points=11]{1}{3}{21}%【同试卷一第三(21)题】


\useproblem[points=11]{1}{3}{22}%【同试卷一第三(22)题】


\useproblem[points=11]{1}{3}{23}%【同试卷一第三(23)题】


\end{document}
